% create a new section called Introduction
\section{Introduction}

\LaTeX is a document preparation system and markup language widely used for the communication and publication of scientific documents in many fields, including mathematics, computer science, engineering, physics, chemistry, and economics. It provides a high-quality typesetting system that is particularly well-suited for documents containing complex mathematical equations and symbols.

\subsection{Easy to learn, easy to use}

\LaTeX is quite easy to learn and use, if you don't need a local installation the best option is clearly to use \Overleaf, a free online latex editor. It integrates a great project manager, allowing to share it to other user and provide a vast range of templates and tutorials for get started with \LaTeX.

For you to learn how to write a document in \LaTeX, a quick introduction is provided in this document. However, if you really want to be able to do everything you want with \LaTeX, there is no secret: you'll have to pass hours to search online, read questions on stackexchange, read the documentation of the packages you want to use, and so on. But don't worry, it's really worth it, and you'll be able to do things that are impossible to do with word.

\subsection{Why use \LaTeX?}

There is several reasons why Latex is widely used in academia and scientific publishing:
\begin{itemize}
  \item \textbf{Open-source and free:} \LaTeX is an open-source software, which means that it is freely available to anyone who wants to use it. 
  \item \textbf{High-quality typesetting:} \LaTeX produces high-quality documents with a professional look and feel. It is particularly well-suited for documents containing complex mathematical equations and symbols, as it provides a powerful set of tools for typesetting mathematics.
  \item \textbf{Customizability:} \LaTeX is highly customizable, allowing users to create their own document classes, packages, and styles. This makes it possible to create documents that are tailored to specific needs and requirements.
  \item \textbf{Cross-platform compatibility:} \LaTeX is available on multiple platforms, including Windows, Mac, and Linux. This makes it easy to collaborate with others who may be using different operating systems.
  \item \textbf{Large community and support:} \LaTeX has a large and active community of users and developers who contribute to the development of the software and provide support to other users. This means that there are many resources available online, including tutorials, forums, and documentation, to help users learn and use \LaTeX effectively.
  \item \textbf{Integration with other tools:} \LaTeX can be easily integrated with other tools and software, such as reference managers, version control systems, and collaborative writing platforms. This makes it a powerful tool for academic writing and publishing.
\end{itemize}

However, to be perfectly honest, one of the used \LaTeX editor and compiler, Overleaf, is not open-source, and a lot of its features are prenium.

\subsection{Overleaf}

Overleaf is a freenium online LaTeX editor that allows users to create, edit, and collaborate on LaTeX documents in real-time. It provides a clean UI for writing and editing LaTeX code, as well as a project manager. However, it's still an online tool, meaning that you need an internet connection to use it, and that you are dependent on the servers of Overleaf.
